\poczatekwykladu{8}

\begin{proof}[Dalszy ciąg dowodu twierdzenia~\ref{tw:dzialania-na-mierzalnych}]
\emph{Suma.} Ustalmy $a\in\R$. Z gęstości liczb wymiernych wynika,
że $f(x)>a-g(x)$ wtedy i tylko wtedy, gdy istnieje $q\in\mathbb{Q}$
takie, że $a-g(x)<q<f(x)$. Zatem
\begin{align*}
  \{x\in X:f(x)+g(x)>a\}
    &=\bigcup_{q\in\mathbb{Q}}
      \bigl(\{x\in X:f(x)>q\}\cap\{x\in X:g(x)>a-q\}\bigr).
\end{align*}
Zbiory w nawiasach należą do $\mathcal{F}$, a suma jest przeliczalna,
więc $f+g$ jest mierzalna.

\emph{Różnica.} Funkcja $-g=(-1)g$ jest mierzalna na mocy
udowodnionej części twierdzenia. Stąd $f-g=f+(-g)$ jest mierzalna.

\emph{Iloczyn.} Korzystamy z tożsamości
\[
  fg=\frac{(f+g)^2-f^2-g^2}{2}.
\]
Wystarczy więc wykazać, że kwadrat funkcji mierzalnej jest mierzalny.
Dla dowolnej takiej funkcji $h$ mamy
\[
  \{x\in X:h(x)^2>a\}
  =\begin{cases}
    X, & a<0,\\
    \{x:h(x)>\sqrt a\}\cup\{x:h(x)<-\sqrt a\}, & a\geq0.
  \end{cases}
\]
W obu przypadkach otrzymujemy zbiór mierzalny.

\emph{Iloraz.} Jego dziedzina
\[
  D=\{x\in X:g(x)\neq0\}
   =\{x:g(x)>0\}\cup\{x:g(x)<0\}
\]
jest mierzalna. Wystarczy sprawdzić mierzalność $1/g$ na $D$,
a następnie zastosować wynik o iloczynie do $f\cdot(1/g)$.
Ponieważ funkcja $t\mapsto1/t$ jest malejąca na każdej z półprostych,
dla $a\in\R$ otrzymujemy
\[
  \{x\in D:1/g(x)>a\}
  =\begin{cases}
    \{x:0<g(x)<1/a\}, & a>0,\\
    \{x:g(x)>0\}, & a=0,\\
    \{x:g(x)>0\}\cup\{x:g(x)<1/a\}, & a<0.
  \end{cases}
\]
Każdy z tych zbiorów należy do $\mathcal{F}|_D$.
\end{proof}

\section{Kresy i granice funkcji mierzalnych}
\label{sec:kresy-granice-mierzalnych}

\begin{twierdzenie}
\label{tw:kresy-granice-mierzalnych}
Niech $(f_n)$ będzie ciągiem funkcji mierzalnych
$f_n\colon X\to\Rbar$. Wtedy funkcje
\[
  \min_{1\leq n\leq N}f_n,\qquad
  \max_{1\leq n\leq N}f_n,\qquad
  \inf_{n\in\N}f_n,\qquad
  \sup_{n\in\N}f_n
\]
są mierzalne. Mierzalne są również
\[
  \liminf_{n\to\infty}f_n,\qquad
  \limsup_{n\to\infty}f_n
\]
oraz granica punktowa $\lim_{n\to\infty}f_n$, jeżeli istnieje
w każdym punkcie $X$.
Wszystkie działania rozumiemy punktowo, na przykład
\[
  \left(\inf_{n\in\N}f_n\right)(x)=\inf_{n\in\N}f_n(x).
\]
\end{twierdzenie}

\begin{proof}
Dla $a\in\R$ mamy
\[
  \left\{x:\min_{1\leq n\leq N}f_n(x)<a\right\}
  =\bigcup_{n=1}^{N}\{x:f_n(x)<a\}\in\mathcal{F}.
\]
Tak samo
\[
  \left\{x:\inf_{n\in\N}f_n(x)<a\right\}
  =\bigcup_{n=1}^{\infty}\{x:f_n(x)<a\}\in\mathcal{F},
\]
co dowodzi mierzalności infimum. Z tożsamości
\[
  \sup_{n\in\N}f_n=-\inf_{n\in\N}(-f_n)
\]
wynika mierzalność supremum; analogicznie otrzymujemy wynik dla maksimum
skończenie wielu funkcji.
Wreszcie korzystamy z punktowych wzorów
\begin{align*}
  \liminf_{n\to\infty}f_n
    &=\sup_{N\in\N}\left(\inf_{n\geq N}f_n\right),\\
  \limsup_{n\to\infty}f_n
    &=\inf_{N\in\N}\left(\sup_{n\geq N}f_n\right).
\end{align*}
Obie funkcje są mierzalne jako kolejne kresy funkcji mierzalnych.
Jeżeli granica istnieje, jest równa zarówno granicy dolnej, jak i górnej,
więc także jest mierzalna.
\end{proof}

\begin{figure}[htbp]
  \centering
  \includegraphics[width=0.74\textwidth]{funkcje-minimum-maksimum.pdf}
  \caption{Punktowe minimum i maksimum dwóch funkcji.}
  \label{fig:minimum-maksimum}
\end{figure}

\section{Funkcje borelowskie i złożenia}
\label{sec:borelowskie-zlozenia}

\begin{definicja}[Funkcja borelowska]
\label{def:funkcja-borelowska}
Funkcję $f\colon\R\to\R$ nazywamy borelowską, jeżeli
dla każdego $a\in\R$ zbiór $\{x\in\R:f(x)>a\}$ jest borelowski.
Jest to mierzalność względem sigma-ciała $\mathcal{B}(\R)$.
\end{definicja}

\begin{uwaga}
W definicji można równoważnie użyć dowolnej z nierówności
$\geq$, $<$ lub $\leq$. Na mocy
twierdzenia~\ref{tw:mierzalnosc-borelowska} warunek ten jest równoważny
temu, że $f^{-1}(B)\in\mathcal{B}(\R)$ dla każdego
$B\in\mathcal{B}(\R)$.
\end{uwaga}

\begin{uwaga}
Każda funkcja borelowska jest mierzalna względem $\mathcal{L}$,
ponieważ $\mathcal{B}(\R)\subseteq\mathcal{L}$.
Mówiąc o funkcji $f\colon\R\to\R$ po prostu „mierzalna”,
będziemy w tym kontekście mieć na myśli mierzalność względem
sigma-ciała zbiorów mierzalnych w sensie Lebesgue'a.
\end{uwaga}

\begin{twierdzenie}[Mierzalność złożenia]
\label{tw:zlozenie-borelowska-mierzalna}
Jeżeli $g\colon X\to\R$ jest mierzalna, a $f\colon\R\to\R$
jest borelowska, to złożenie
\[
  f\circ g\colon X\to\R,\qquad (f\circ g)(x)=f(g(x))
\]
jest mierzalne.
\end{twierdzenie}

\begin{proof}
Dla każdego $a\in\R$ zachodzi
\[
  \{x:f(g(x))>a\}
  =(f\circ g)^{-1}((a,+\infty))
  =g^{-1}\bigl(f^{-1}((a,+\infty))\bigr).
\]
Przeciwobraz $f^{-1}((a,+\infty))$ jest borelowski, więc jego
przeciwobraz przez mierzalną funkcję $g$ należy do $\mathcal{F}$.
\end{proof}

\begin{twierdzenie}[Funkcje ciągłe]
\label{tw:ciagla-borelowska}
Każda funkcja ciągła $f\colon\R\to\R$ jest borelowska.
\end{twierdzenie}

\begin{proof}
Dla każdego $a\in\R$ przeciwobraz otwartej półprostej
$(a,+\infty)$ przez funkcję ciągłą jest otwarty, a więc borelowski.
Schemat takiego przeciwobrazu przedstawia
rysunek~\ref{fig:ciagla-przeciwobraz}.
\end{proof}

\begin{figure}[htbp]
  \centering
  \includegraphics[width=0.72\textwidth]{funkcja-ciagla-przeciwobraz.pdf}
  \caption{Przeciwobraz półprostej $(a,+\infty)$ przez funkcję ciągłą.}
  \label{fig:ciagla-przeciwobraz}
\end{figure}

\section{Równość prawie wszędzie}
\label{sec:prawie-wszedzie}

Niech $(X,\mathcal{F},\mu)$ będzie przestrzenią z miarą,
a $f,g\colon X\to\R$ będą mierzalne.
Poniższe pojęcia stosujemy również do funkcji o wartościach w $\Rbar$.

\begin{definicja}[Równość prawie wszędzie]
\label{def:rownosc-prawie-wszedzie}
Mówimy, że $f$ i $g$ są równe $\mu$-prawie wszędzie, jeżeli
\[
  \mu(\{x\in X:f(x)\neq g(x)\})=0.
\]
Piszemy wtedy $f=g$ $\mu$-p.w. Używa się także skrótu
$\mu$-a.e. (\emph{almost everywhere}).
\end{definicja}

Zbiory występujące w tej definicji są mierzalne. Również dla wartości
rozszerzonych można bowiem zapisać
\[
  \{x:f(x)<g(x)\}
  =\bigcup_{q\in\mathbb{Q}}
    \bigl(\{x:f(x)<q\}\cap\{x:g(x)>q\}\bigr),
\]
a zbiór $\{f\neq g\}$ jest sumą zbiorów $\{f<g\}$ i $\{g<f\}$.

\begin{przyklad}[Funkcja Dirichleta]
\label{prz:dirichlet-prawie-wszedzie}
Na $\R$ rozważmy funkcje
\[
  O(x)=0,\qquad D(x)=\mathbf{1}_{\mathbb{Q}}(x).
\]
Zbiór, na którym się różnią, jest równy $\mathbb{Q}$.
Jest on przeliczalny, zatem ma miarę Lebesgue'a zero i
$O=D$ $\lambda$-p.w. Funkcje te nie są jednak równe $\delta_0$-p.w.,
ponieważ $\delta_0(\mathbb{Q})=1$.
Z kolei $D$ i funkcja stale równa $1$ są równe $\delta_0$-p.w.:
ich zbiór różnic to $\R\setminus\mathbb{Q}$, który nie zawiera zera.
\end{przyklad}

\begin{twierdzenie}
\label{tw:prawie-wszedzie-rownowaznosc}
Równość $\mu$-prawie wszędzie jest relacją równoważności
na zbiorze funkcji mierzalnych.
\end{twierdzenie}

\begin{proof}
Zwrotność i symetria wynikają bezpośrednio z definicji.
Dla przechodniości załóżmy, że $f=g$ $\mu$-p.w. oraz
$g=h$ $\mu$-p.w. Rozbijamy zbiór różnic na dwie rozłączne części:
\begin{align*}
  \mu(\{f\neq h\})
    &=\mu(\{f\neq h\}\cap\{f=g\})
      +\mu(\{f\neq h\}\cap\{f\neq g\})\\
    &\leq\mu(\{g\neq h\})+\mu(\{f\neq g\})=0.
\end{align*}
W pierwszym składniku użyliśmy tego, że z $f(x)=g(x)$ i
$f(x)\neq h(x)$ wynika $g(x)\neq h(x)$.
Zatem $f=h$ $\mu$-p.w.
\end{proof}

\begin{definicja}[Własność zachodząca prawie wszędzie]
\label{def:wlasnosc-prawie-wszedzie}
Mówimy, że własność $P(x)$ zachodzi $\mu$-prawie wszędzie,
jeżeli istnieje zbiór $E\in\mathcal{F}$ taki, że
$\mu(E^c)=0$ i $P(x)$ jest prawdziwa dla każdego $x\in E$.
\end{definicja}
